\begin{abstract}

 % In this paper, we introduce \ours, a We introduce the task of scientific fact-checking.  Given a corpus of scientific articles and a claim about a scientific finding, a fact-checking model must identify abstracts that support or refute the claim. In addition, it must provide \emph{rationales} for its predictions in the form of evidentiary sentences from the retrieved abstracts. For this task, we introduce \ours, a dataset of 1.4K expert-written scientific claims paired with evidence-containing abstracts annotated with labels and rationales. We develop a pipeline system which retrieves potentially relevant abstracts, identifies rationale sentences from those abstracts, and makes final \supports / \refutes decisions based on the content of the selected rationales. and assess its performance on \ours. We observe , while fact-checking models trained on Wikipedia articles or political news have difficulty generalizing to our task, simple domain adaptation techniques represent a promising avenue for improvement. Finally, we provide initial results showing how our model can be used to verify claims relevant to COVID-19 on the CORD-19 corpus.

  % Uncomment URL for camera-ready.
  % Our dataset will be made publicly available at \website.


We introduce scientific claim verification, a new task to select abstracts from the research literature containing evidence that supports or refutes a given scientific claim, and to identify rationales justifying each decision. To study this task, we construct \ours, a dataset of 1.4K expert-written scientific claims paired with evidence-containing abstracts annotated with labels and rationales. We develop baseline models for \ours, and demonstrate that these models benefit from combined training on a large dataset of claims about Wikipedia articles, together with the new \ours data. We show that our claim verification system is able to identify plausible evidence for 23 / 36 claims relevant to COVID-19 on the CORD-19 corpus. Our results and experiments strongly suggest that our new task and data will support significant future research efforts.\footnote{Data, code, and a web demo will released publicly.}


\end{abstract}